\chapter{Ekuacionet jolineare, sistemet lineare dhe vlerat vetjake}
\begin{qellime}
Studenti zbaton përgjysmimin, pikën fikse, Njutonin dhe sekanten; zgjidh sisteme lineare me eliminim e faktorizim; interpreton kushtëzimin; llogarit vlera e vektorë vetjakë me metodën fuqi dhe algoritmin QR.
\end{qellime}

\section{Ekuacioni skalar $f(x)=0$}
Rrënja është pikë ku modeli plotëson një kusht fizik: ekuilibër force, energji e caktuar, kohë goditjeje ose parametër kalibrimi. Metoda duhet zgjedhur sipas informacionit të disponueshëm, intervalit të sigurt dhe kostos së vlerësimit të funksionit.

\subsection{Përgjysmimi}
Nëse $f$ është e vazhdueshme dhe $f(a)f(b)<0$, ekziston të paktën një rrënjë në $[a,b]$. Pas $n$ hapash,
\begin{equation}
 |x_n-x_*|\le \frac{b-a}{2^{n+1}}.
\end{equation}
Metoda është e ngadaltë, por robuste. Ajo nuk dallon rrënjët me shumëfish çift, sepse shenja mund të mos ndryshojë.

\begin{shembull}[Algoritmi i përgjysmimit]
Jepen $f$, $a$, $b$ dhe toleranca $\tau$, me $f(a)f(b)<0$. Për sa kohë $(b-a)/2>\tau$, vendoset $c=(a+b)/2$. Nëse $f(a)f(c)\le0$, zëvendësohet $b$ me $c$; përndryshe zëvendësohet $a$ me $c$. Në fund kthehet $(a+b)/2$.
\end{shembull}

\subsection{Pika fikse dhe Njutoni}
Shkrimi $x=g(x)$ jep iteracionin $x_{n+1}=g(x_n)$. Konvergjenca lokale kërkon $|g'(x_*)|<1$. Përzgjedhja e $g$ është pjesë e metodës; forma algjebrike e papërshtatshme mund të divergojë.

Linearizimi i $f$ rreth $x_n$ jep metodën e Njutonit,
\begin{equation}
 x_{n+1}=x_n-\frac{f(x_n)}{f'(x_n)}.
\end{equation}
Pranë një rrënje të thjeshtë, gabimi është kuadratik: $|e_{n+1}|\approx C|e_n|^2$. Megjithatë, derivati i vogël, pika fillestare e keqe ose dalja nga domeni mund ta dështojnë metodën. Një strategji praktike kombinon intervalin e përgjysmimit me hapin e Njutonit.

Sekantja zëvendëson derivatin me pjerrësinë ndërmjet dy pikave. Ajo nuk kërkon derivat dhe ka rend konvergjence $\varphi=(1+\sqrt5)/2$.

\begin{figure}[ht]
\centering\includegraphics[width=.75\textwidth]{03_konvergjenca_rrenjeve.pdf}
\caption{Konvergjenca lineare e përgjysmimit dhe konvergjenca e shpejtë e Njutonit për $\cos x-x=0$.}
\end{figure}

\section{Sistemet jolineare}
Për $\vect F(\vect x)=\vect0$, linearizimi përdor Jakobin $J_{ij}=\partial F_i/\partial x_j$:
\begin{equation}
 J(\vect x_k)\Delta\vect x_k=-\vect F(\vect x_k),\qquad
 \vect x_{k+1}=\vect x_k+\Delta\vect x_k.
\end{equation}
Në çdo hap zgjidhet një sistem linear; inversi nuk formohet. Shkallëzimi i ndryshoreve është i rëndësishëm kur komponentët kanë njësi ose madhësi shumë të ndryshme. Në probleme të vështira përdoret amortizimi $\vect x_{k+1}=\vect x_k+\lambda\Delta\vect x_k$, $0<\lambda\le1$.

\begin{shembull}[Ekuilibri mekanik]
Për dy masa të lidhura me susta jolineare, $\vect F(\vect x)$ përmbledh forcat neto. Jakobiani është matrica e ngurtësisë tangjente. Konvergjenca e Njutonit tregon jo vetëm zgjidhjen, por edhe afërsinë ndaj humbjes së stabilitetit kur $J$ bëhet pothuaj singular.
\end{shembull}

\section{Sistemet lineare}
Eliminimi i Gausit e shndërron $A$ në formë trekëndore. Pivotimi i pjesshëm zgjedh elementin me modul më të madh në kolonën aktive dhe shmang pjesëtimin me numra të vegjël. Faktorizimi shkruhet
\begin{equation}
 PA=LU,
\end{equation}
ku $P$ është matricë permutimi, $L$ trekëndore e poshtme dhe $U$ trekëndore e sipërme. Për shumë anë të djathta, faktorizimi kryhet një herë dhe zëvendësimet përsëriten.

Për matrica simetrike pozitivisht të përcaktuara përdoret Cholesky,
\begin{equation}
 A=LL^T,
\end{equation}
që shfrytëzon simetrinë dhe ka rreth gjysmën e kostos së LU. Kostoja e faktorizimit të një matrice të dendur $n\times n$ është $\Oh(n^3)$; ruajtja është $\Oh(n^2)$.

\begin{kujdes}
Determinanti i vogël nuk është masë e mirë e kushtëzimit, sepse varet nga shkalla. Përdoret $\kappa(A)$ ose vlera singulare më e vogël. Gjithashtu, një mbetje e vogël $\vect r=\vect b-A\tilde{\vect x}$ nuk garanton gabim të vogël kur $A$ është e keqkushtëzuar.
\end{kujdes}

\section{Vlerat dhe vektorët vetjakë}
Problemi
\begin{equation}
 A\vect v=\lambda\vect v
\end{equation}
shfaqet në modet normale, stabilitet, mekanikë kuantike dhe analiza të rrjeteve. Për një sistem masë--sustë,
\begin{equation}
 K\vect u=\omega^2 M\vect u.
\end{equation}
Vektorët vetjakë janë format modale; vlerat vetjake japin frekuencat katrore.

\begin{figure}[ht]
\centering\includegraphics[width=.72\textwidth]{05_modet_normale.pdf}
\caption{Mode normale të një sistemi njëpërmasor me skaje të fiksuara.}
\end{figure}

\subsection{Metoda fuqi}
Iteracioni $\vect x_{k+1}=A\vect x_k/\|A\vect x_k\|$ izolon vektorin e lidhur me vlerën vetjake me modul më të madh, nëse komponenti fillestar përkatës nuk është zero. Shpejtësia përcaktohet nga $|\lambda_2/\lambda_1|$. Vlera vlerësohet me kuocientin Rayleigh
\begin{equation}
 \rho(\vect x)=\frac{\vect x^T A\vect x}{\vect x^T\vect x}.
\end{equation}

\subsection{Algoritmi QR}
Faktorizimi $A_k=Q_kR_k$ dhe përditësimi $A_{k+1}=R_kQ_k$ prodhojnë matrica të ngjashme, pra me të njëjtat vlera vetjake. Për matrica simetrike, $A_k$ afrohet drejt formës diagonale. Në praktikë përdoren reduktimi Hessenberg/tridiagonal dhe zhvendosjet për konvergjencë të shpejtë.

\section{Renditja e faqeve si problem vetjak}
Një matricë kalimi $P$ përshkruan probabilitetin e kalimit ndërmjet faqeve. Vektori stacionar plotëson $\vect p=P\vect p$. Faktori i amortizimit shmang nyjet pa dalje dhe komponentët e shkëputur:
\begin{equation}
 G=\alpha P+(1-\alpha)\frac{\vect1\vect1^T}{n}.
\end{equation}
Metoda fuqi mbi $G$ jep renditjen. Ky shembull lidh algjebrën lineare, zinxhirët Markov dhe simulimin.

\begin{lidhje}
Laboratori ndërton sistemin masë--sustë, verifikon simetrinë, llogarit frekuencat dhe vizaton modet. Një ushtrim shtesë ndërton një rrjet të vogël faqesh dhe krahason iteracionin e fuqisë me zgjidhjen e drejtpërdrejtë.
\end{lidhje}

\section*{Pyetje kontrolli}
\begin{enumerate}
\item Kur preferohet përgjysmimi ndaj Njutonit?
\item Pse pivotimi është pjesë e qëndrueshmërisë së eliminimit?
\item Çfarë përcakton shpejtësinë e metodës fuqi?
\end{enumerate}
\section{Algoritmet kryesore dhe zbatimi në Python}
Ky kapitull përmban tri familje problemesh. Secila kërkon një kriter ndalimi të lidhur me mbetjen, jo vetëm me ndryshimin ndërmjet dy iteracioneve.

\subsection{Përgjysmimi dhe Njutoni i mbrojtur}
\begin{lstlisting}[language=Python,caption={Përgjysmimi me kontroll të intervalit.}]
def pergjysmimi(f, a, b, tol=1e-12, max_iter=200):
    fa, fb = f(a), f(b)
    if fa * fb > 0:
        raise ValueError("Intervali nuk kufizon nje ndryshim shenje.")
    for _ in range(max_iter):
        c = 0.5 * (a + b)
        fc = f(c)
        if abs(fc) < tol or 0.5*(b-a) < tol:
            return c
        if fa * fc <= 0:
            b, fb = c, fc
        else:
            a, fa = c, fc
    raise RuntimeError("Nuk u arrit konvergjenca.")
\end{lstlisting}

\begin{lstlisting}[language=Python,caption={Njutoni me kontroll të derivatit.}]
def njutoni(f, df, x, tol=1e-12, max_iter=50):
    for _ in range(max_iter):
        fx, dfx = f(x), df(x)
        if abs(fx) < tol:
            return x
        if abs(dfx) < np.sqrt(np.finfo(float).eps):
            raise ArithmeticError("Derivat shume i vogel.")
        hapi = -fx / dfx
        x = x + hapi
        if abs(hapi) < tol * (1 + abs(x)):
            return x
    raise RuntimeError("Nuk u arrit konvergjenca.")
\end{lstlisting}
Në një zbatim robust, hapi i Njutonit pranohet vetëm nëse qëndron brenda intervalit dhe zvogëlon mbetjen; përndryshe kryhet një përgjysmim.

\subsection{Zgjidhja e sistemit linear}
\begin{lstlisting}[language=Python,caption={Zgjidhja dhe diagnostika pa formuar inversin.}]
A = np.array([[4., -1., 0.], [-1., 4., -1.], [0., -1., 3.]])
b = np.array([15., 10., 10.])
x = np.linalg.solve(A, b)
r = b - A @ x
print("x =", x)
print("mbetja relative =", np.linalg.norm(r)/np.linalg.norm(b))
print("kappa_2 =", np.linalg.cond(A))
\end{lstlisting}
Operatori \texttt{solve} përdor faktorizime të përshtatshme. Shkrimi \texttt{inv(A) @ b} shmanget.

\subsection{Metoda fuqi}
\begin{lstlisting}[language=Python,caption={Vlera vetjake dominante dhe kuocienti Rayleigh.}]
def metoda_fuqi(A, tol=1e-11, max_iter=10_000):
    x = np.ones(A.shape[0])
    x /= np.linalg.norm(x)
    lam_old = 0.0
    for k in range(max_iter):
        y = A @ x
        x = y / np.linalg.norm(y)
        lam = x @ A @ x
        if abs(lam - lam_old) < tol * (1 + abs(lam)):
            return lam, x, k + 1
        lam_old = lam
    raise RuntimeError("Metoda fuqi nuk konvergoi.")
\end{lstlisting}
Kontrolli përfundimtar është norma $\|A\vect x-\lambda\vect x\|$. Shenja e vektorit vetjak nuk ka kuptim fizik më vete: $\vect x$ dhe $-\vect x$ përfaqësojnë të njëjtin mod.

\subsection{Lidhja me modet normale}
Për $K\vect u=\omega^2M\vect u$, përdoret \texttt{scipy.linalg.eigh(K, M)} kur $K$ dhe $M$ janë simetrike dhe $M$ pozitivisht e përcaktuar. Frekuencat renditen, normalizohen dhe kontrollohet ortogonaliteti $U^TMU=I$.

Rendi i paraqitjes së metodave të rrënjëve, sistemeve dhe vlerave vetjake ndjek traditën e teksteve shqip \cite{hanelli,hoxha,datja}; termat e algjebrës lineare numerike kontrollohen edhe kundrejt punimeve të Departamentit të Matematikës së Aplikuar të UT-së \cite{zeqiri2020}.

\subsection{Konvergjenca e metodave për rrënjët}
Le të jetë $x_*$ rrënjë e thjeshtë. Në iteracionin e pikës fikse $x_{k+1}=g(x_k)$, teorema e vlerës mesatare jep
\begin{equation}
 e_{k+1}=x_{k+1}-x_*=g'(\xi_k)(x_k-x_*)=g'(\xi_k)e_k.
\end{equation}
Nëse $|g'(x)|\le q<1$ në një interval që ruhet nga $g$, atëherë $|e_k|\le q^k|e_0|$. Ky është konvergjim linear. Për Njutonin, zhvillimi Taylor rreth $x_k$ jep
\begin{equation}
 0=f(x_*)=f(x_k)+f'(x_k)(x_*-x_k)+\frac{f''(\xi_k)}{2}(x_*-x_k)^2.
\end{equation}
Pas zëvendësimit të hapit të Njutonit merret
\begin{equation}
 e_{k+1}=\frac{f''(\xi_k)}{2f'(x_k)}e_k^2.
\end{equation}
Konvergjenca kuadratike vlen vetëm pasi iteracioni ka hyrë në fqinjësinë e rrënjës dhe kur $f'(x_*)\ne0$. Për një rrënjë me shumëfish $m$, Njutoni standard bëhet linear; formula e modifikuar $x_{k+1}=x_k-mf(x_k)/f'(x_k)$ rikthen rendin kuadratik kur $m$ dihet.

Metoda e Müller-it ndërton një polinom kuadratik në tri pika dhe zgjedh rrënjën e tij që vazhdon trajektoren e iteracionit. Ajo mund të prodhojë rrënjë komplekse edhe nga pika fillestare reale, çka është avantazh për polinomet, por kërkon trajtim të kujdesshëm të degës së rrënjës katrore dhe të anulimit në emërues.

\subsection{Eliminimi i Gausit dhe pivotimi}
Në hapin $k$, elementi $a_{kk}^{(k)}$ përdoret si pivot dhe shumëzuesit janë
\begin{equation}
 \ell_{ik}=\frac{a_{ik}^{(k)}}{a_{kk}^{(k)}},\qquad i=k+1,\ldots,n.
\end{equation}
Rreshti $i$ përditësohet me $a_{ij}^{(k+1)}=a_{ij}^{(k)}-\ell_{ik}a_{kj}^{(k)}$. Në aritmetikë të saktë, kjo është një sekuencë transformimesh elementare. Në aritmetikë me saktësi të fundme, pivot i vogël prodhon shumëzues të mëdhenj dhe zmadhon gabimet. Pivotimi i pjesshëm zgjedh
\begin{equation}
 p=\arg\max_{i\ge k}|a_{ik}^{(k)}|
\end{equation}
dhe këmben rreshtat $p$ dhe $k$. Rezultati ruhet si $PA=LU$. Faktorizimi kushton afërsisht $2n^3/3$ veprime me presje të lëvizshme, kurse dy zëvendësimet për çdo anë të djathtë kushtojnë $\Oh(n^2)$.

Formimi i $A^{-1}$ kërkon zgjidhjen e $n$ sistemeve dhe pastaj një shumëzim tjetër me $\vect b$; ai shton punë dhe rrumbullakime pa nevojë. Për këtë arsye \texttt{solve(A,b)} është zgjedhja standarde.

\subsection{Njutoni për sisteme dhe kriteret e ndalimit}
Nga zhvillimi shumëpërmasor
\begin{equation}
 \vect F(\vect x+\Delta\vect x)=\vect F(\vect x)+J(\vect x)\Delta\vect x+\Oh(\|\Delta\vect x\|^2)
\end{equation}
vendosja e anës së majtë në zero jep sistemin e Njutonit. Duhet të kontrollohen njëkohësisht norma e mbetjes dhe madhësia relative e hapit,
\begin{equation}
 \|\vect F(\vect x_k)\|\le\tau_F,
 \qquad
 \|\Delta\vect x_k\|\le\tau_x(1+\|\vect x_k\|).
\end{equation}
Vetëm kriteri i hapit mund të pranojë stagnim; vetëm mbetja mund të mashtrojë kur ekuacionet kanë shkallë shumë të ndryshme. Shkallëzimi i secilit ekuacion me një madhësi karakteristike i bën kriteret fizikisht të krahasueshme.

\begin{lstlisting}[language=Python,caption={Njutoni i amortizuar për një sistem jolinear.}]
import numpy as np

def njutoni_sistem(F, J, x0, tol=1e-10, maksimumi=50):
    x = np.array(x0, dtype=float)
    for k in range(maksimumi):
        Fx = np.asarray(F(x), dtype=float)
        if np.linalg.norm(Fx, ord=np.inf) <= tol:
            return x, k
        hapi = np.linalg.solve(J(x), -Fx)
        alfa = 1.0
        # Kërkim prapa: pranohet vetëm ulja e normës së mbetjes.
        while np.linalg.norm(F(x + alfa*hapi)) >= np.linalg.norm(Fx):
            alfa *= 0.5
            if alfa < 2.0**-20:
                raise RuntimeError("Njutoni stagnoi; kontrolloni shkallëzimin.")
        x = x + alfa*hapi
    raise RuntimeError("U arrit numri maksimal i iteracioneve.")
\end{lstlisting}

\subsection{Metoda fuqi dhe algoritmi QR: derivim spektral}
Nëse $A$ ka një bazë vektorësh vetjakë dhe $\vect x_0=\sum_i c_i\vect v_i$, atëherë
\begin{equation}
 A^k\vect x_0=\lambda_1^k\left[c_1\vect v_1+
 \sum_{i=2}^{n}c_i\left(\frac{\lambda_i}{\lambda_1}\right)^k\vect v_i\right].
\end{equation}
Kur $|\lambda_1|>|\lambda_2|$ dhe $c_1\ne0$, termat e tjerë shuhen relativisht si $|\lambda_2/\lambda_1|^k$. Normalizimi shmang tejmbushjen, por nuk ndryshon drejtimin.

Në algoritmin QR, $A_k=Q_kR_k$ dhe $A_{k+1}=R_kQ_k=Q_k^TA_kQ_k$, kështu që $A_{k+1}$ është e ngjashme me $A_k$. Për matrica simetrike, transformimet ortogonale ruajnë mirë normën dhe diagonalja konvergon te vlerat vetjake. Zhvendosja $A_k-\mu_kI=Q_kR_k$, e ndjekur nga $A_{k+1}=R_kQ_k+\mu_kI$, përshpejton ndarjen e vlerave të afërta.

\subsection{Përmbledhje dhe ushtrime}
Metodat e këtij kapitulli bashkohen nga linearizimi: Njutoni linearizon funksionin, eliminimi zgjidh problemin linear dhe analiza spektrale shpjegon mënyrën si vepron matrica në drejtime të veçanta.
\begin{enumerate}
\item Nxirrni rendin e sekantes duke supozuar $|e_{k+1}|\simeq C|e_k e_{k-1}|$.
\item Implementoni Müller-in me numra kompleksë dhe provojeni te $x^4+1=0$.
\item Numëroni veprimet e eliminimit të Gausit dhe nxirrni termin $2n^3/3$.
\item Për një sistem masë--sustë, verifikoni ortogonalitetin e modeve ndaj matricës së masës.
\item Ndërtoni një rrjet faqesh dhe studioni ndikimin e parametrit të amortizimit në konvergjencën e renditjes.
\end{enumerate}

\subsection{Metoda e Müller-it dhe rrënjët komplekse}
Metoda e Müller-it përdor tri pika $x_0,x_1,x_2$ dhe ndërton polinomin kuadratik interpolues për $f$. Duke vendosur $h_1=x_1-x_0$, $h_2=x_2-x_1$ dhe diferencat e pjesëtuara
\begin{equation}
 \delta_1=\frac{f(x_1)-f(x_0)}{h_1},\qquad
 \delta_2=\frac{f(x_2)-f(x_1)}{h_2},
\end{equation}
koeficienti kuadratik është $d=(\delta_2-\delta_1)/(h_1+h_2)$. Polinomi rreth $x_2$ shkruhet
\begin{equation}
 p(x_2+s)=f(x_2)+bs+ds^2,qquad b=\delta_2+h_2d.
\end{equation}
Rrënja që jep hapin më të qëndrueshëm është
\begin{equation}
 s=-\frac{2f(x_2)}{b+\operatorname{sgn}_*(b)\sqrt{b^2-4df(x_2)}},
\end{equation}
ku shenja zgjidhet që emëruesi të ketë modul maksimal. Kjo zgjedhje shmang anulimin ndërmjet $b$ dhe rrënjës katrore. Pika e re është $x_3=x_2+s$, pastaj ruhet tresha më e fundit.

Diskriminanti mund të jetë negativ edhe kur pikat janë reale, prandaj implementimi duhet të lejojë aritmetikë komplekse. Kjo e bën Müller-in të dobishëm për polinome pa rrënjë reale. Renditja e konvergjencës është afërsisht $1.84$, më e lartë se e sekantes, por çdo hap përdor më shumë informacion dhe mund të jetë më pak i parashikueshëm. Për probleme ku kërkohet një rrënjë reale brenda një intervali fizik, një metodë e mbrojtur me përgjysmim është zakonisht më e sigurt.

\subsubsection{Refinimi iterativ i sistemeve lineare}
Pas zgjidhjes së $A\widehat x=b$, llogaritet mbetja $r=b-A\widehat x$ dhe zgjidhet
\begin{equation}
 A\delta x=r,qquad \widehat x\leftarrow\widehat x+\delta x.
\end{equation}
Në aritmetikë të saktë, një hap do të korrigjonte të gjithë gabimin. Në praktikë, nëse mbetja llogaritet me saktësi më të lartë se faktorizimi, refinimi mund të rikuperojë disa shifra. Faktorizimi LU ripërdoret, kështu që çdo korrigjim kushton vetëm dy zëvendësime trekëndore. Procesi ndalet kur mbetja e shkallëzuar nuk ulet më ose korrigjimi bëhet i papërfillshëm.

Një kriter i dobishëm komponentor është
\begin{equation}
 \eta=\max_i\frac{|r_i|}{(|A||\widehat x|+|b|)_i}.
\end{equation}
Ai interpretohet si perturbimi relativ minimal i afërt i të dhënave që e bën $\widehat x$ zgjidhje të saktë. Një $\eta$ pranë saktësisë së makinës tregon algoritëm të qëndrueshëm në prapavijë, por gabimi përpara mund të mbetet i madh kur sistemi është i keqkushtëzuar.

\subsubsection{Problemi i përgjithësuar i vlerave vetjake}
Në mekanikën e sistemeve me shumë shkallë lirie shfaqet
\begin{equation}
 K\vect u=\omega^2M\vect u,
\end{equation}
ku $M$ është matrica e masës dhe $K$ matrica e ngurtësisë. Nëse $M$ është simetrike pozitivisht e përcaktuar, faktorizimi Cholesky $M=LL^T$ e kthen problemin në
\begin{equation}
 L^{-1}KL^{-T}\vect z=\omega^2\vect z,qquad \vect u=L^{-T}\vect z.
\end{equation}
Matrica e transformuar është simetrike. Modet mund të normalizohen që $\vect u_i^TM\vect u_j=\delta_{ij}$ dhe atëherë $\vect u_i^TK\vect u_j=\omega_i^2\delta_{ij}$. Këto identitete janë prova numerike: mbetja $\|K\vect u_i-\omega_i^2M\vect u_i\|$ kontrollon çdo çift vetjak, ndërsa ortogonaliteti kontrollon të gjithë bazën.

Në rrjete ose sisteme me simetri, vlera vetjake të shumëfishta nuk përcaktojnë vektorë unikë; përcaktohet vetëm nënhapësira vetjake. Krahasimi komponent më komponent i dy vektorëve nga programe të ndryshme mund të jetë i pakuptimtë. Duhet krahasuar mbetja, vlera vetjake dhe hapësira e shtrirë nga modet përkatëse.

